% =====================================================================
% 4 Metodologia do Ensaio 1
% Fontes: paper/secoes/04-identificacao.tex, 05-dados.tex;
% docs/pre-especificacao.md; docs/ars/17 (resposta ao parecer).
% Posição já corrigida pelo parecer de 22/09/2026 — ver pauta do orientador.
% =====================================================================

\chapter{Metodologia do Ensaio 1}
\label{cap:metodologia}

A exposição segue a ordem parâmetro-alvo, hipóteses de identificação e
estimação. As escolhas descritas neste capítulo são as do plano de análise
fechado em 21 de setembro de 2026, com os desvios posteriores registrados no
Apêndice~\ref{ap:preespecificacao}.

\section{Parâmetro-alvo}
\label{sec:estimando}

A fonte de variação é a intensidade de exposição anterior à lei entre os 184
municípios cearenses. O tratamento é contínuo, e o estimador é o de
\citet{callaway2024continuous}, que separa dois parâmetros:

\begin{alineas}
  \item $\ATT$, o efeito médio, entre os municípios com dose $d$, de ter
    recebido $d$ em vez de zero. É um efeito no nível da dose e é identificado
    sob tendências paralelas no resultado não tratado;
  \item $ATE(d)$ e sua derivada $ACR(d)$, a curva dose-resposta causal. A curva
    exige tendências paralelas fortes (SPT), que restringem trajetórias sob
    doses que o município não recebeu e excluem seleção por ganhos.
\end{alineas}

O plano de análise declarou a curva como alvo primário. A estimação mostrou que
a derivada agregada não é interpretável (Seção~\ref{sec:res-nivel}), e desde 22
de setembro de 2026 o resultado principal é o $\ATT$. O desvio está registrado
com data e motivo.

A variável de dose é área plantada de banana, e não aplicação aérea. O que o
desenho estima é, portanto, o \textbf{efeito de forma reduzida da proibição
segundo uma \emph{proxy} de intensidade da bananicultura}. Ele não é, sem
hipóteses adicionais, o efeito da retirada da pulverização aérea onde ela
ocorria. \decisao{1}{declarar o estimando como forma reduzida segundo a
\emph{proxy}; a alternativa, dose-resposta da pulverização aérea, exigiria o
cadastro operacional anterior à lei.}

Sobre a hipótese forte, o teste de tendências anteriores não a isola, porque a
implicação testável antes do tratamento é comum às duas versões do paralelismo.
Ele pode, porém, falsificá-la: \citet{callaway2024continuous} propõem verificar
se a relação entre a variação do desfecho e a dose é nula no período anterior,
e o projeto implementa essa verificação (Seção~\ref{sec:res-robustez}).

\section{Variável de tratamento}
\label{sec:dose}

A dose é a área plantada média de 2015 a 2018 da cultura-âncora, na Produção
Agrícola Municipal do IBGE, normalizada. Usa-se área plantada, e não colhida,
porque o período de 2012 a 2017 foi de seca severa no semiárido e a área
colhida responde à quebra de safra.

A cultura-âncora foi escolhida pelo dado, e não por suposição. O melão tem área
positiva em 10 municípios e o algodão em 28, suporte insuficiente para a
estimação não paramétrica. A banana tem área positiva em 169 municípios, com
índice de Gini de 0,86 e 78\% da área no decil superior. É também a cultura que
corresponde à química documentada: a procimidona é fungicida de bananal.

Três limitações de medida acompanham a escolha.

\begin{alineas}
  \item A proibição é de método. Área plantada mede intensidade agrícola, e
    culturas e municípios diferem em quanto recebem aplicação por aeronave. A
    Seção~\ref{sec:diluicao} mostra que essa distância é grande.
  \item A janela de medida começa depois da notícia da proibição. O projeto de
    lei foi apresentado em 24 de fevereiro de 2015, e a janela de 2015 a 2018 é
    posterior a esse marco. Se produtores responderam à expectativa, a dose já
    incorpora parte do ajuste. \decisao{3}{manter 2015--2018 como janela
    principal e reportar 2010--2014 como sensibilidade; promover 2010--2014 a
    principal muda o plano de análise e pode mudar o grupo tratado.}
  \item O zero da PAM é zero de \emph{proxy}. O § 2º do art. 28-B alcança o
    controle vetorial por aeronave, de modo que um município sem banana pode
    ter sido tratado (Seção~\ref{sec:comparacao}).
\end{alineas}

\section{Exposição gestacional}
\label{sec:exposicao}

Um desfecho perinatal não é contemporâneo à exposição. Um nascimento de março
de 2019 foi gestado, em dois terços, antes da vigência da lei. O painel calcula,
para cada célula município-mês, a fração da janela gestacional posterior à
proibição. Apenas as coortes nascidas a partir de outubro de 2019 têm gestação
integralmente posterior. A janela gestacional é fixa: usar a duração observada
condicionaria a exposição a um desfecho do estudo, a prematuridade. As coortes
cuja gestação atravessa a vigência são descartadas na estimação principal, o
que corresponde a 9,4\% das células. As séries de intoxicação aguda usam o mês
do evento, sem retroprojeção.

\section{Grupo de comparação}
\label{sec:comparacao}

O grupo de comparação é o dos municípios de dose nula, construído de quatro
maneiras: área nula da cultura-âncora; área nula de qualquer cultura candidata
a aplicação aérea; ausência de operador aeroagrícola registrado; e baixa
aptidão agroclimática. A definição primária é a quarta, e as demais entram como
sensibilidade. O grupo tem 15 municípios (14 na quarta definição).

Que esse grupo seja nunca tratado é hipótese, e não fato estabelecido. O Censo
Agropecuário de 2006 a apoia para a aplicação agrícola: nenhum município de
dose nula tinha estabelecimento com aplicação por aeronave. O Censo não mede o
controle vetorial aéreo, e não há fonte que o meça com suficiência (o SINAN
registra 89 notificações por agente de uso em saúde pública em oito anos). O
peso dessa hipótese é alto, porque os mesmos 15 municípios servem de controle a
todos os estimadores binários do projeto.

Duas verificações relativas ao grupo tratado foram concluídas. A varredura das
leis municipais anteriores a 2019 nos 17 municípios do decil superior encontrou
uma única proibição, a de Limoeiro do Norte, revogada em 2010; não há,
portanto, município já tratado na janela de 2015 a 2018. E o índice de aptidão
agroclimática FAO-GAEZ, construído para os 184 municípios pela receita de
\citet{reynier2025glyphosate}, é maior onde há mais banana (média de 0,37 no
decil superior contra 0,26 nos demais). \decisao{2}{tratar a aptidão como
construção do grupo de comparação e checagem de relevância; estimação por
variável instrumental fica como extensão declarada, porque a aptidão no
semiárido se correlaciona com disponibilidade de água, que afeta o desfecho
diretamente.}

\section{Dados}
\label{sec:dados}

A Tabela~\ref{tab:fontes} resume as fontes. O painel do Ensaio 1 é município
por mês, balanceado, com 184 municípios e 96 meses (2015 a 2022), num total de
17.664 células, das quais 17.654 têm nascimento registrado. Os desfechos são o
peso ao nascer, o baixo peso (menos de 2.500 g) e a prematuridade (menos de 37
semanas), agregados à célula. A Tabela~\ref{tab:descritivas} traz as
descritivas. \decisao{4}{estender a série do SINASC até 19/12/2024, que já
está disponível na fonte original, refazer painel e análise de poder e
congelar a data de extração de cada base.}

\begin{table}[htb]
\centering
\caption{Fontes de dados e estado de aquisição}
\label{tab:fontes}
\small
\begin{tabular}{p{3.6cm}p{2.1cm}p{8.0cm}}
\toprule
Fonte & Estado & Uso e observação \\
\midrule
PAM/IBGE         & adquirida   & Dose: área plantada por cultura, 2015--2018 \\
SINASC/DATASUS   & adquirida   & Desfechos de nascimento, 2015--2022 \\
SIM/DATASUS      & adquirida   & Óbito fetal \\
SINAN            & adquirida   & Intoxicação exógena; 1.217 notificações por
                                 agrotóxico de uso agrícola \\
População (IBGE) & adquirida   & Denominador; 2022 vem do Censo \\
FAO-GAEZ         & adquirida   & Aptidão agroclimática, 184 municípios \\
Censo Agropecuário 2006 & adquirida & Estabelecimentos por equipamento de
                                 aplicação, incluída a aeronave \\
Leis municipais  & levantadas  & 17 municípios do decil superior \\
SISAGUA          & suspensa    & Quebra de registro laboratorial no Ceará a
                                 partir de 2020 \\
ANA (bacias)     & suspensa    & Depende de coordenada que o SISAGUA não traz \\
INMET            & suspensa    & Direção do vento sem variação espacial \\
MapBiomas        & descartada  & Banana em classe genérica \\
SEMACE, ADAGRI, MAPA & a solicitar & Registros de aplicação aérea anteriores
                                 à lei, por pedido de acesso à informação \\
\bottomrule
\end{tabular}
\legend{Fonte: elaboração própria. Estado em setembro de 2026.}
\end{table}

\begin{table}[htb]
\centering
\caption{Descritivas do painel, Ceará, 2015--2022}
\label{tab:descritivas}
\small
\begin{tabular}{lr}
\toprule
Grandeza & Valor \\
\midrule
Nascidos vivos (após limpeza)                 & 1.001.709 \\
Peso médio ao nascer                          & 3.208 g \\
Baixo peso ($<$ 2.500 g)                      & 8,34\% \\
Prematuridade ($<$ 37 semanas)                & 12,58\% \\
Óbitos fetais                                 & 11.240 \\
Notificações por agrotóxico agrícola (SINAN)  & 1.217 \\
Municípios com dose positiva                  & 169 \\
Municípios com dose nula                      & 15 \\
\bottomrule
\end{tabular}
\legend{Fonte: elaboração própria, a partir de SINASC, SIM, SINAN e PAM.}
\end{table}

Três fontes foram obtidas, inspecionadas e suspensas, e o motivo de cada
suspensão é informação sobre o desenho. No SISAGUA, três das quatro moléculas
mais frequentes na Chapada não constam da lista monitorada, e as detecções de
agrotóxico no Ceará caem de 44 em 2019 para zero em 2020, 2021 e 2022, com mais
amostras coletadas; um modelo de diferenças em diferenças sobre essa variável
atribuiria à lei uma mudança de prática laboratorial. Sem o SISAGUA, o desenho
de montante e jusante fica sem insumo. No INMET, a coerência de direção do
vento entre estações é de 0,96: os alísios sopram de leste, e estar a favor do
vento equivale a estar a oeste da fonte. O canal de intoxicação, previsto com
dados de internação, passou ao SINAN, porque a série hospitalar registra um
único caso de intoxicação acidental por agrotóxico no período.

\section{Estimação e inferência}
\label{sec:estimacao}

A estimação principal usa o estimador de \citet{callaway2024continuous} com
aproximação não paramétrica, sobre o painel colapsado em um período anterior e
um posterior. A implementação impõe três restrições que afetam o que se pode
afirmar: opera sobre dois períodos; não produz o estudo de eventos e a curva
pelo mesmo procedimento; e não acomoda covariáveis.

O efeito agregado no nível é, na implementação usada, a diferença entre a
variação média do desfecho nos 169 municípios de dose positiva e a variação
média nos 15 de dose nula. Na inferência, portanto, o grupo escasso é o de
controle. Reportam-se, para o nível, o erro-padrão do \emph{bootstrap}
multiplicador, o \emph{wild bootstrap} com hipótese nula imposta
\cite{cameron2008wild}, uma análise de permutação e a exclusão sucessiva de
cada município de controle. A permutação da dose é reportada como análise de
permutação, e não como inferência exata: a exatidão exigiria que a dose tivesse
sido atribuída de forma intercambiável, o que não ocorre, porque a área de
banana segue aptidão e geografia. \decisao{5}{adotar como procedimento principal
para o nível uma inferência que respeite os 15 controles, na linha de
\citet{conley2011} com os papéis de tratado e controle invertidos.}

A camada de robustez compreende o poder do teste de tendências anteriores
\cite{roth2023trending}, a sensibilidade a violações de paralelismo
\cite{rambachan2023credible}, a verificação da inclinação em dose em datas
fictícias do período anterior e o controle sintético com diferenças em
diferenças \cite{arkhangelsky2021sdid}, na versão intraestadual. Uma
especificação com pareamento por escore de propensão e diferenças em diferenças
\cite{bravoureta2020psmdid} entra como camada de comunicação dos resultados.

\section{Poder, falsificação e multiplicidade}
\label{sec:poder}

O desenho declara seu poder antes de estimar. O efeito mínimo detectável
previsto para o peso ao nascer era de 33,4 g. O critério de falsificação fixa
um piso de efeito esperado de 15 g, abaixo do extremo inferior do intervalo de
23 a 32 g de \citet{reynier2025glyphosate}, por se supor atenuação quando se
remove um método e não uma molécula. Se a meia-largura do intervalo de
confiança ficar acima do piso, o resultado é reportado como limite informativo
sobre o parâmetro de forma reduzida, e não como achado.

A família confirmatória tem duas hipóteses, peso médio e óbito fetal, com
correção de Holm. Todo o restante é exploratório e declarado como tal. O plano
veda trocar o desfecho ou descer de unidade geográfica depois de ver o
resultado.

\section{Ameaças à identificação}
\label{sec:ameacas}

O Quadro~\ref{qua:ameacas} resume as ameaças e o tratamento dado a cada uma.

\begin{quadro}[htb]
\centering
\caption{Ameaças à identificação e resposta do desenho}
\label{qua:ameacas}
\small
\begin{tabular}{|p{3.6cm}|p{4.9cm}|p{5.6cm}|}
\hline
\textbf{Ameaça} & \textbf{Efeito sobre o resultado} & \textbf{Resposta} \\
\hline
Erro de classificação do tratamento
  & Atenua o efeito estimado; é o mecanismo da falta de poder
  & Medir a aplicação aérea por registro administrativo; concentrar a
    comparação onde havia aeronave \\
\hline
Prática interrompida antes de 2019
  & Se a pulverização cessou na Chapada em 2011, a lei não removeu exposição
  & Consultar a ação civil pública de 2011 e os relatórios do MAPA \\
\hline
Grupo de comparação exposto
  & Desloca o nível; pode inverter o sinal
  & Quatro construções do grupo; exclusão sucessiva dos controles \\
\hline
Antecipação desde 02/2015
  & Contamina a dose, não só o desfecho
  & Janela de dose 2010--2014 como sensibilidade; estudo de eventos até 2018 \\
\hline
Tendências paralelas fortes
  & A curva deixa de ser identificada
  & Resultado principal no nível; verificação em datas fictícias \\
\hline
Seleção para nascimento vivo
  & Efeito real pode aparecer nulo ou invertido
  & Óbito fetal como segundo desfecho confirmatório \\
\hline
Fiscalização desconhecida
  & O parâmetro é de intenção de tratar
  & Pedidos de acesso à informação; canal de intoxicação \\
\hline
Poucos grupos
  & Precisão aparente maior que a real
  & Inferência que respeita os 15 controles \\
\hline
\end{tabular}
\legend{Fonte: elaboração própria.}
\end{quadro}
